\documentclass{article}
\usepackage[UTF8]{ctex}
\usepackage{amsmath,amssymb,amsfonts}
\usepackage{enumitem}
\usepackage{geometry}
\geometry{a4paper, margin=1in}
% === 错题本排版优化（不参与题目标识） ===
\setlength{\parindent}{0pt}
\setlength{\parskip}{0.35em plus 0.10em minus 0.08em}
\setlength{\emergencystretch}{3em}
\setlist{itemsep=0.25em, topsep=0.40em, parsep=0pt}
\allowdisplaybreaks[2]
% === 错题本排版优化结束 ===
\begin{document}
\section*{中国矿业大学2026年高等代数真题}
\section*{21．中国矿业大学2026年高等代数试卷}
一．填空题．每题5分，共50分．\\

\subsection{题目1}
1．多项式 \(x^{3}+4 x^{2}+7 x+12\) 的有理根为 \(\_\_\_\_\)．

\subsection{题目2}
2．设 \(A\) 是 3 阶实方阵，且 \(|E+A|=0,|E+2 A|=0,|E-3 A|=0, E\) 为 3 阶单位矩阵，则 \(|E+6 A|=\)\\
\(\_\_\_\_\)．

\subsection{题目3}
3．已知 \(\alpha_{1}, \alpha_{2}, \alpha_{3}, \alpha_{4}\) 为4维列向量组，且 \(\alpha_{1}, \alpha_{2}, \alpha_{4}\) 线性无关，\(\alpha_{3}=\alpha_{1}+\alpha_{2}-\alpha_{4}, \beta=\alpha_{1}+\alpha_{3}+3 \alpha_{4}\)，则非齐次线性方程组 \(A X=\beta\) 的通解为 \(\_\_\_\_\)．

\subsection{题目4}
4．设 \(\alpha, \beta\) 为3维列向量，且 \(\alpha \beta^{\mathrm{T}}=\left(\begin{array}{ccc}\frac{1}{2} & -\frac{1}{2} & \frac{1}{2} \\ -\frac{1}{3} & \frac{1}{3} & -\frac{1}{3} \\ \frac{1}{6} & -\frac{1}{6} & \frac{1}{6}\end{array}\right)\) ，则 \(\beta^{\mathrm{T}} \alpha=\) \(\_\_\_\_\)

\subsection{题目5}
5．若 \(A\) 是 5 阶实可逆矩阵，则 \(A\) 与 \(-5 A\) \(\_\_\_\_\)（填＂可能＂或＂不可能＂）合同。\\
\(\_\_\_\_\)

\subsection{题目6}
6．设 \(n\) 阶矩阵 \(A=-E_{11}-2 E_{22}-3 E_{33}-\cdots-n E_{n n}\) ，则与 \(A\) 可交换的矩阵所组成的线性空间的维数为 \(\_\_\_\_\)．

\subsection{题目7}
7．若二次型 \(f\left(x_{1}, x_{2}, x_{3}\right)=2 x_{1}^{2}+2 x_{2}^{2}+3 x_{3}^{2}+\lambda x_{1} x_{2}+2 x_{2} x_{3}+2 x_{1} x_{3}\) 是正定的，则 \(\lambda\) 的取值范围是\\
\(\_\_\_\_\)。

\subsection{题目8}
8．定义线性变换 \(\sigma\left(x_{1}, x_{2}, x_{3}\right)=\left(x_{1}+3 x_{2}+x_{3}, x_{2}+x_{3}, x_{1}-x_{3}\right)\) ，则 \(\sigma\) 的秩为 \(\_\_\_\_\)．

\subsection{题目9}
9．设 \(A\) 是一个秩为3的四阶方阵，\(A\) 的对角元的代数余子式分别为1，－2，3，－4，则 \(A\) 的伴随矩阵\(A^{*}\) 的特征值为 \(\_\_\_\_\)。

\subsection{题目10}
10．已知 \(A\) 是奇数阶正交矩阵，且 \(|A|=1\) ，则 1 \(\_\_\_\_\)（填＂一定＂或＂不一定＂）是 \(A\) 的特征值．

二．（10 分）设 \(f_{1}(x), f_{2}(x), \cdots, f_{s}(x), g_{1}(x), g_{2}(x), \cdots, g_{t}(x)\) 均为多项式，证明：\(f_{1}(x) f_{2}(x) \cdots f_{s}(x)\)与 \(g_{1}(x) g_{2}(x) \cdots g_{t}(x)\) 互素的充要条件是对任意的 \(1 \leq i \leq s, 1 \leq j \leq t\) 都有 \(f_{i}(x)\) 与 \(g_{j}(x)\) 互素．

三。（10 分）已知 \(A\) 为 \(n\) 阶实矩阵，\(A^{\mathrm{T}}\) 为 \(A\) 的转置，\(\beta\) 为 \(n\) 维实列向量。

\begin{itemize}
  \item[（1）] （5 分）证明：齐次线性方程组 \(A^{\mathrm{T}} A X=0\) 与 \(A X=0\) 同解。
  \item[（2）] （5 分）证明：方程组 \(A^{\mathrm{T}} A X=A^{\mathrm{T}} \beta\) 必然有解．
\end{itemize}

四．（15 分）设 \(A, B\) 为 \(n\) 阶实方阵，\(A^{2}=A, B^{2}=B\) ，且 \(E-A-B\) 可逆，证明：\(A, B\) 的秩相同．

五．（15 分）设 \(A\) 是 \(n\) 阶正定矩阵，\(B\) 为 \(n\) 阶实反称矩阵，证明：\(A-B^{2}\) 可逆．

六．（15 分）设 \(A, B, C, D\) 是数域 \(P\) 上两两可交换的 \(n\) 阶方阵，且 \(A C+B D=E\) ，记\(V=\left\{X \in P^{n} \mid A B X=0\right\}, V_{1}=\left\{X \in P^{n} \mid A X=0\right\}, V_{2}=\left\{X \in P^{n} \mid B X=0\right\}\) ．证明：\(V=V_{1} \oplus V_{2}\) ．

七．（15 分）设 \(\sigma\) 是 \(n\) 维欧氏空间 \(V\) 上的正交变换，记

\[
V_{1}=\{\alpha \in V \mid \sigma(\alpha)=\alpha\}, V_{2}=\{\alpha-\sigma(\alpha) \mid \alpha \in V\} .
\]

证明：\(V=V_{1} \oplus V_{2}\) ．

八．（20 分）已知 \(A\) 是 \(n(n \geq 3)\) 阶矩阵，\(A^{*}\) 为 \(A\) 的伴随矩阵，且 \(A^{*}=A^{2}-2 A-E,|A|=-2\) 。

\begin{itemize}
  \item[（1）] （5 分）证明：\(A\) 可相似对角化．
  \item[（2）] （12 分）分别求 \(n=3,4,5,6\) 时，\(A\) 的相似对角形．
  \item[（3）] （3 分）求 \(A\) 的迹（用 \(n\) 表示）．
\end{itemize}
\end{document}
